\section{2. Related Work}

\subsection{Traditional Light Field Depth Estimation Methods}

Many depth estimation methods for light-field cameras have been proposed \cite{lightfielddepthestimationE}. Among the earliest and most influential approaches, Wanner et al. \cite{Wannergloballyconsistentdep} proposed a variational framework that derives depth from the slope of linear structures in Epipolar Plane Images (EPI). By formulating depth estimation as an energy minimization problem, they explicitly enforce photo-consistency across angular views, where the disparity $d$ is analytically derived from the EPI slope. This approach achieves significantly accurate disparity maps in controlled environments. However, most of these EPI-based methods rely heavily on the global Lambertian assumption and work poorly on glossy or specular surfaces \cite{Lambertianassumptionglossys}. When the photo-consistency constraint is violated by view-dependent reflections, the linear EPI structures become distorted or fragmented, causing the depth prediction to completely fail or degrade into mere texture copying.

While EPI-based methods primarily exploit angular correspondence, other approaches leverage focal stack cues to address the under-constrained problem in texture-less or non-Lambertian regions. Tao et al. \cite{Taodepthfromcombiningdefoc} proposed a clustering method that eliminates specular pixels when enforcing photo-consistency, combining defocus and correspondence cues to estimate depth. Specifically, they attempt a binary classification of pixels into either Lambertian or specular, utilizing the defocus cue to recover depth for the rejected specular regions. Although this strategy successfully mitigates the impact of strong highlights, it cannot handle general glossy surfaces with spatially-varying BRDFs (SVBRDFs) \cite{spatiallyvaryingBRDFsglossy}. Furthermore, such heuristic pixel rejection or binary masking lacks physical interpretability and struggles with complex mixed materials where diffuse and specular components are intricately blended.

To tackle the limitations of heuristic masking, recent work has explored integrating physical reflectance models into light-field depth estimation. Johannsen et al. \cite{Johannsendepthestimationand} formulated the depth recovery as a joint optimization problem by incorporating the dichromatic model to separate diffuse and specular components. By explicitly modeling the specular lobe, they demonstrated improved depth accuracy on objects with uniform dielectric materials. Similarly, Kim et al. \cite{KimdepthestimationofnonLa} leveraged a simplified BRDF to regularize the cost volume, attempting to disentangle geometry from reflectance. Nonetheless, these traditional physical BRDF modeling approaches suffer from extremely high computational complexity due to the iterative non-convex optimization required to solve for both depth and reflectance parameters simultaneously. Moreover, they are typically limited to simplified, single-lobe BRDFs and fail to properly generalize to real-world scenes characterized by complex mixed materials and diverse illumination conditions.

In short, traditional light-field depth estimation methods exhibit three fundamental limitations. First, they severely rely on the global Lambertian assumption; in non-Lambertian (specular, scattering) regions, the violation of view-consistency leads to a complete failure of depth prediction, often degenerating into texture copying. Second, traditional physical BRDF modeling entails prohibitively high computational complexity and struggles to effectively handle complex mixed materials ubiquitous in real-world scenarios. Third, these methods lack an in-depth analysis of the angular sampling frequency characteristics of light fields, rendering them incapable of distinguishing complex reflection types from limited angular samples in a lightweight and efficient manner. In contrast, our proposed framework substantially resolves these bottlenecks. Specifically, our MRI-inspired angular frequency analysis breaks through the limitation of distinguishing complex reflections from limited samples, providing a physically meaningful and reliable prior for subsequent depth estimation. Moreover, the dual-mask physical modeling unifies the representation of Lambertian, non-Lambertian, and mixed scenes from a physical optics perspective, endowing the data-driven model with explicit physical interpretability to adaptively handle mixed materials. Finally, the component-aware adaptive fusion strategy comprehensively eliminates the fatal defect of traditional EPI methods in non-Lambertian regions, significantly improving depth accuracy and robustness in complex mixed scenes.

\subsection{Deep Learning-based General Light Field Depth Estimation}

To address the limitations of hand-crafted features in traditional geometry-based methods, deep learning-based approaches have been extensively explored for general light field depth estimation. Shin et al. \cite{EPINetsubapertureimages} proposed EPINet, a pioneering fully convolutional network that estimates depth from sub-aperture images (SAIs). By extracting spatial and angular features from multi-directional SAIs and fusing them via a multi-stream architecture, EPINet successfully captures the parallax cues embedded in the 4D light field. Subsequent variants, such as EPINet4Dir \cite{EPINet4DirdirectionalSAIs}, further refine this paradigm by explicitly modeling directional disparities to enhance angular correspondence. While these SAI-based convolutional neural networks (CNNs) demonstrate remarkable accuracy on standard Lambertian datasets, they primarily rely on implicit feature matching. Consequently, when the photo-consistency assumption is violated by view-dependent reflections, the learned spatial-angular correlations become unreliable, causing the depth prediction to struggle with glossy or specular surfaces.

While SAI-based methods focus on 2D spatial and 1D angular slices, more recent approaches attempt to directly process the 4D light field tensor to capture global angular correlations. For instance, LFattNet \cite{LFattNet4Dconvolutionattent} leverages 4D convolutions and attention mechanisms to model the complex spatial-angular dependencies across the entire light field, significantly improving depth estimation in texture-less regions. Furthermore, Wang et al. \cite{DistgSSRdisentangledfeatures} proposed DistgSSR, which introduces a disentangled spatial-angular feature extraction module. By decoupling the spatial and angular dimensions and employing multi-scale feature fusion, DistgSSR effectively mitigates the aliasing bimodal distribution problem often encountered in cost volume-based methods. These advanced baselines substantially enhance the representational capacity of the network. However, they still formulate depth estimation as a purely data-driven regression task, treating the 4D light field merely as a high-dimensional signal without considering the underlying physical image formation process.

Despite the remarkable representational capacity of these deep learning-based models, they share several fundamental limitations when generalizing to complex real-world environments. First, they are predominantly pure data-driven black-box models that lack physical interpretability at the optical level, making it difficult to explicitly model the light-matter interaction process. Second, these networks typically treat the highlights or scattering in non-Lambertian regions as mere noise or ordinary textures, failing to explicitly distinguish material reflection types, which leads to a dramatic performance drop in Non-Lambertian and Mixed scenes. Furthermore, the unified network architecture cannot adaptively adjust the depth estimation strategy for different reflection regions, causing the error gradients from non-Lambertian areas to interfere with the overall optimization. In contrast to these purely empirical approaches, our work addresses these issues by introducing a Dual-Mask Physical Modeling framework that explicitly unifies the formulation of Lambertian, Non-Lambertian, and Mixed scenes from a physical optics perspective. By incorporating an MRI-inspired angular frequency analysis for material reflection parsing and a component-aware adaptive fusion strategy, our model substantially resolves the fatal flaw of previous methods where predictions degenerate into texture copying in non-Lambertian regions, thereby achieving notable robustness in mixed-material scenarios.

\subsection{Deep Learning-based Non-Lambertian and Material-Aware Depth Estimation}

While the aforementioned SAI- and 4D-based methods primarily focus on implicit feature matching under the Lambertian assumption, recent efforts have explicitly incorporated material awareness and dual cues to tackle non-Lambertian scenes. MaterialDualCueNet \cite{MaterialDualCueNetdualcuesd} proposes a multi-stream architecture that jointly exploits defocus and correspondence cues. Specifically, it introduces an auxiliary material classification branch to identify glossy regions and adaptively weight the depth features. While this approach significantly improves depth accuracy on isolated specular objects, it heavily relies on a complex learning-based classifier and additional hardware cues (e.g., the defocus branch). Our preliminary experiments indicate that such data-driven material classifiers struggle to learn effective joint depth and material representations, particularly when ground-truth annotations for diverse spatially-varying BRDFs (SVBRDFs) are scarce.

In contrast to joint estimation, another line of research attempts to circumvent the violation of photo-consistency by explicitly separating reflection components prior to depth regression. For instance, Dichromatic-LFNet \cite{DichromaticLFNetdichromatic} leverages the dichromatic model to decompose light field images into diffuse and specular components. By purifying the diffuse albedo, it successfully restores the photo-consistency assumption for subsequent Epipolar Plane Image (EPI) based depth estimation. However, this strategy substantially treats specular reflections as noise to be eliminated. Consequently, the depth estimation policy fails to thoroughly decouple and adaptively fuse reflection components  [i.e., diffuse, specular, and scattering]. In complex mixed scenes, discarding the specular component inevitably leads to severe texture loss and sub-optimal precision, as the geometric cues embedded in the specular lobes are entirely ignored.

To address the physical inconsistency of purely data-driven separation, more recent approaches have explored physics-informed priors to model non-Lambertian reflections. Physics-Informed LF Depth \cite{PhysicsInformedLFDepthphys} formulates the depth estimation as an under-constrained problem and incorporates angular frequency analysis to distinguish between Lambertian and non-Lambertian regions. By applying physical constraints on the epipolar plane, it demonstrates enhanced robustness on translucent and glossy surfaces. Nonetheless, this method fails to unify the mathematical modeling of Lambertian and non-Lambertian reflections from a fundamental physical optics perspective  [e.g., angular frequency domain and wavevector deflection]. Specifically, it lacks a fine-grained dual-mask physical constraint tailored for both the medium properties and angular directions, limiting its generalization to arbitrary mixed materials.

In short, while existing deep learning-based non-Lambertian and material-aware methods have made significant strides, they share critical limitations. First, current material classification predominantly relies on complex learning-based classifiers or extra hardware cues (e.g., defocus branches), which struggle to extract effective depth and material features. Second, they fail to unify the mathematical modeling of Lambertian and non-Lambertian surfaces from a physical optics perspective, lacking fine-grained dual-mask physical constraints for medium and angular directions. Finally, their depth estimation strategies do not thoroughly decouple and adaptively fuse reflection components, resulting in sub-optimal performance in complex mixed scenes. To address these challenges, our proposed Unified Dual-Mask Physical Model departs from conventional data-driven paradigms. We introduce an MRI-inspired angular frequency analysis for physically meaningful material parsing, formulate a dual-mask physical modeling to unify diverse reflection types, and design a component-aware adaptive fusion strategy, thereby dramatically enhancing depth estimation robustness in non-Lambertian and mixed scenarios.

Despite their progress, existing methods either struggle with strict Lambertian assumptions or rely on purely data-driven material classifiers lacking physical interpretability, which often leads to severe depth degradation (e.g., texture copying) when the view-consistency hypothesis is violated in complex mixed-reflectance scenes. Inspired by these observations, we propose a unified dual-mask physical modeling approach that addresses the aforementioned limitations by integrating an MRI-inspired angular frequency analysis for physically meaningful material parsing, a physically interpretable dual-mask mechanism to uniformly model diverse reflectance types, and a component-aware adaptive fusion strategy to robustly prevent depth degradation in non-Lambertian regions.